\clearpage\thispagestyle{empty}
\pdfbookmark[0]{Quellen zur Berechnung des Pessachfestes (1802)}{osterfest1802-sources}
\begin{center}\textbf{Quellen zur Berechnung des Pessachfestes (1802)}\end{center}
\noindent Diese KI-Übertragung des deutschen Originals in heutiges Deutsch folgt Gauß’ \textit{Berechnung des jüdischen Osterfestes} in den \textit{Werken}, Band~VI (Göttingen, 1874), S.~80--81. Der Erstdruck erschien in Franz Xaver von Zachs \textit{Monatlicher Correspondenz zur Beförderung der Erd- und Himmels-Kunde}, Band~5, Mai 1802, Aufsatz~XLVI, S.~435--437, unter dem Namen Dr.~Gauß in Braunschweig. \href{https://www.digitale-sammlungen.de/de/view/bsb10538597?page=447}{Erstdruck}. \href{https://gdz.sub.uni-goettingen.de/id/PPN236007467}{Historischer deutscher Text in den Werken}.
\par\smallskip\noindent Titel und Text geben Gauß’ Ausdruck \textit{jüdisches Osterfest} als \textit{Pessachfest} wieder. Die Kalenderbegriffe und Zahlenformen der \textit{Werke} bleiben erhalten, einschließlich der über das Monatsende hinausgezählten Märztage im Beispiel. Zum heutigen Sprachgebrauch siehe Barbara Staudinger, \href{https://www.jmw.at/news/die_pessach-formel}{\textit{Die Pessach-Formel}}, Jüdisches Museum Wien, 27.~März 2025.
\par\smallskip\noindent Übertragung und Satz wurden mit KI erstellt.
\clearpage
